% ===========================================================================
%  Annexe — Recueil de sujets
% ===========================================================================
\chapter{Sujets de baccalauréat}
\markboth{Sujets}{Sujets}

Cette annexe rassemble ce qu'il faut savoir de l'épreuve elle-même : sa
structure, un sujet complet à traiter dans les conditions réelles, et un
état des lieux des sessions disponibles pour la série OSE au moment de la
rédaction de ce livre. Les exercices repris dans les chapitres proviennent
tous de ce sujet ; on trouve ici de quoi les replacer dans leur contexte.

% ---------------------------------------------------------------------------
\section*{La structure de l'épreuve}

L'épreuve de mathématiques générales de la série OSE dure quatre heures et
pèse le coefficient cinq. Le code matière est le $009$, et la calculatrice
scientifique non programmable est autorisée — les mêmes conditions que pour
les séries C et D, mais un coefficient inférieur à celui de la série S.

La session 2026, seule disponible en intégralité au moment de la rédaction
de ce livre, dessine une composition en quatre blocs :

\begin{center}
\setlength{\tabcolsep}{6pt}
\renewcommand{\arraystretch}{1.35}
\begin{tabular}{@{}L{1.8cm}L{1.3cm}L{5.8cm}L{2.4cm}@{}}
\toprule
\textsf{\footnotesize\bfseries PARTIE} &
\textsf{\footnotesize\bfseries POINTS} &
\textsf{\footnotesize\bfseries CONTENU} &
\textsf{\footnotesize\bfseries CHAPITRES} \\
\midrule
Exercice 1 & 4 & probabilités et variable aléatoire & 13, 14 \\
Exercice 2 & 7 (Ex. 2+3) & statistique à deux variables & 15 \\
Exercice 3 & 7 (Ex. 2+3) & mathématiques financières & annexe \\
Problème & 9 & analyse : fonction, tangente, réciproque, suite & 1 à 7 \\
\bottomrule
\end{tabular}
\end{center}

\noindent
Le sujet transcrit ne détaille pas la répartition exacte des sept points
entre les exercices 2 et 3 : les deux exercices sont présentés ici sans
barème détaillé, fidèlement à ce que porte le document source. Ce que l'on
peut affirmer avec certitude, en revanche, est la place que tient le
problème d'analyse : neuf points sur vingt, presque la moitié de l'épreuve,
et un déroulé récurrent — étude d'une fonction associant logarithme et
exponentielle, fonction auxiliaire pour étudier son signe, tangente,
réciproque d'une restriction, puis une suite définie par une intégrale.

Aucune des quatre sessions OSE identifiées — 2021, 2022, 2023, 2026 — ne
comporte de nombres complexes, de barycentre ou de similitudes : ce constat,
déjà fait au moment de composer les chapitres 9 à 12, explique pourquoi ce
livre les traite comme un bloc à part (voir l'annexe de progression).

% ---------------------------------------------------------------------------
\section*{Un sujet complet : session 2026, série OSE}

À traiter en quatre heures, sans interruption, calculatrice scientifique non
programmable autorisée.

\subsection*{Exercice 1 \textnormal{(4 points)}}

Une urne contient trois plaquettes indiscernables au toucher, numérotées
$1$, $2$ et $3$. On tire successivement et avec remise trois plaquettes ; on
note dans l'ordre les trois numéros obtenus. On considère les événements
\[
  A : \text{« au moins une plaquette porte un numéro pair »},
\]
\[
  B : \text{« le produit des trois numéros obtenus vaut $6$ »} .
\]
Soit $X$ la variable aléatoire égale au produit des trois numéros obtenus.

\begin{enumerate}[label=\textbf{\arabic*.}, leftmargin=*, itemsep=0.25em]
  \item Calculer $\proba{A}$ et $\proba{B}$. \hfill\textup{(1 pt)}
  \item Déterminer $X(\Omega)$, puis dresser le tableau de la loi de
        probabilité de $X$. \hfill\textup{(1,25 pt)}
  \item Calculer $\Esp(X)$ et $\Var(X)$. \hfill\textup{(0,75 pt)}
  \item On note $F$ la fonction de répartition de $X$. Représenter
        graphiquement $F$. \hfill\textup{(0,5 pt)}
  \item On répète trois fois, de manière indépendante, l'épreuve consistant à
        tirer les trois plaquettes avec remise. Calculer la probabilité
        d'obtenir exactement deux fois l'événement $B$. \hfill\textup{(0,5 pt)}
\end{enumerate}

\subsection*{Exercice 2 \textnormal{(sur les 7 points des exercices 2 et 3)}}

Le tableau suivant donne la production $y_{i}$ (en tonnes) de maïs d'un
Petit Périmètre Irrigué, observée de 2019 à 2024, $x_{i}$ désignant le rang
de l'année :
\[
  \begin{array}{c|cccccc}
    x_{i} & 1 & 2 & 3 & 4 & 5 & 6 \\
    \hline
    y_{i} & 60 & 65 & 70 & 74 & 75 & 80
  \end{array}
\]
On admet que $\sum x_{i} = 21$, $\sum y_{i} = 424$, $\sum x_{i}^{2} = 91$,
$\sum y_{i}^{2} = 30\,226$ et $\sum x_{i}y_{i} = 1\,551$.
\begin{enumerate}[label=\textbf{\arabic*.}, leftmargin=*, itemsep=0.2em]
  \item Représenter le nuage de points et placer le point moyen $G$.
  \item Calculer le coefficient de corrélation linéaire de la série, et
        interpréter le résultat.
  \item Déterminer, par la méthode des moindres carrés, une équation de la
        droite de régression $(D)$ de $y$ en $x$.
  \item En supposant que l'évolution se poursuive selon le même modèle,
        estimer la production attendue en 2028.
\end{enumerate}

\subsection*{Exercice 3 \textnormal{(sur les 7 points des exercices 2 et 3)}}

\textbf{Partie I.} Un capital $C$ est placé à intérêt simple pendant
$15$ mois ; sa valeur acquise dépasse le capital initial de $10\,\%$.
\begin{enumerate}[label=\textbf{\alph*)}, leftmargin=*, itemsep=0.15em]
  \item Calculer le taux annuel $t$ de ce placement.
  \item Sachant que l'intérêt produit vaut $40\,000$~Ar, déterminer $C$ puis
        la valeur acquise.
\end{enumerate}

\textbf{Partie II.} On place désormais un capital $C = 400\,000$~Ar à
intérêt composé au taux annuel de $6\,\%$, pendant $n$ périodes.
\begin{enumerate}[label=\textbf{\alph*)}, leftmargin=*, itemsep=0.15em]
  \item Déterminer la durée $n$ nécessaire pour que la valeur acquise atteigne
        $535\,290{,}23$~Ar.
  \item En posant $C_{n} = C(1+t)^{n}$, montrer que
        $t = e^{\frac1n\ln\left(\frac{C_{n}}{C}\right)}-1$.
  \item En utilisant cette formule, déterminer le taux annuel nécessaire
        pour que ce même capital atteigne $600\,000$~Ar au bout de $5$~ans.
\end{enumerate}

\textbf{Partie III.} On verse $6$ annuités constantes au taux annuel de
$10\,\%$, et l'on obtient une valeur acquise de $771\,561{,}00$~Ar.
\begin{enumerate}[label=\textbf{\alph*)}, leftmargin=*, itemsep=0.15em]
  \item Déterminer le montant de chaque annuité.
\end{enumerate}
On verse à présent $6$ annuités constantes de $100\,000$~Ar chacune, qui
produisent une valeur acquise de $811\,518{,}90$~Ar.
\begin{enumerate}[label=\textbf{\alph*)}, leftmargin=*, itemsep=0.15em]
  \setcounter{enumii}{1}
  \item Déterminer le taux annuel $t$ de ce second placement.
\end{enumerate}

\subsection*{Problème \textnormal{(9 points)}}

On considère la fonction numérique $f$ définie sur $\left]-1\,;+\infty\right[$
par
\[
  f(x) = e^{-x}\ln(x+1) .
\]
On désigne par $(C)$ sa courbe représentative dans un repère orthonormé
$(O\,;\vecu,\vecv)$ d'unité $2$ cm.

\begin{enumerate}[label=\textbf{\arabic*.}, leftmargin=*, itemsep=0.3em]
  \item Soit $g$ la fonction définie sur $\left]-1\,;+\infty\right[$ par
        $g(x) = \ln(x+1) - \dfrac{1}{x+1}$.
        \begin{enumerate}[label=\textbf{\alph*)}, leftmargin=1.4em]
          \item Étudier le sens de variation de $g$ (on ne demande pas sa
                courbe représentative). \hfill\textup{(0,5 pt)}
          \item Montrer que l'équation $g(x) = 0$ admet une solution unique
                $\alpha$ sur $\left]-1\,;+\infty\right[$ où
                $\dfrac12 < \alpha < 1$. \hfill\textup{(0,5 pt)}
          \item En déduire le signe de $g(x)$ suivant les valeurs de $x$.
                \hfill\textup{(0,5 pt)}
        \end{enumerate}
  \item \begin{enumerate}[label=\textbf{\alph*)}, leftmargin=1.4em]
          \item Calculer $\lim_{x\to-1^{+}}f(x)$. Interpréter géométriquement
                ce résultat. \hfill\textup{(0,5 pt)}
          \item Sachant que $f(x) = \dfrac{x}{e^{x}}\times\dfrac{\ln(x+1)}{x}$,
                calculer $\lim_{x\to+\infty}f(x)$. Conclure.
                \hfill\textup{(0,5 pt)}
        \end{enumerate}
  \item Montrer que, pour tout $x \in \left]-1\,;+\infty\right[$,
        $f'(x) = \dfrac{-g(x)}{e^{x}}$. \hfill\textup{(0,5 pt)}
  \item Dresser le tableau de variation de $f$. \hfill\textup{(0,5 pt)}
  \item \begin{enumerate}[label=\textbf{\alph*)}, leftmargin=1.4em]
          \item Démontrer que $f(\alpha) = \dfrac{e^{-\alpha}}{\alpha+1}$.
                \hfill\textup{(0,5 pt)}
          \item Écrire l'équation de la tangente $(T)$ à $(C)$ au point
                d'abscisse $x_{0} = 0$. \hfill\textup{(0,5 pt)}
        \end{enumerate}
  \item Tracer $(T)$ et $(C)$ dans le même repère (on prend
        $\alpha = 0{,}8$). \hfill\textup{(0,75 pt)}
  \item Soit $h$ la restriction de $f$ à l'intervalle $I = \left]-1\,;0\right]$.
        \begin{enumerate}[label=\textbf{\alph*)}, leftmargin=1.4em]
          \item Montrer que $h$ admet une fonction réciproque $h^{-1}$ sur un
                intervalle $J$ que l'on déterminera. \hfill\textup{(0,75 pt)}
          \item Dresser le tableau de variation de $h^{-1}$ sur $J$.
                \hfill\textup{(0,5 pt)}
          \item Calculer $h(0)$ et $\left(h^{-1}\right)'(0)$.
                \hfill\textup{(0,75 pt)}
          \item Tracer dans le même repère que $(C)$ la courbe représentative
                $(C^{-1})$ de la fonction réciproque $h^{-1}$ de $h$.
                \hfill\textup{(0,5 pt)}
        \end{enumerate}
  \item Soit $\left(U_{n}\right)_{n\in\N}$ la suite définie par
        $U_{n} = \displaystyle\int_{a_{n}}^{a_{n+1}}\dfrac{\ln(x+1)}{x+1}\dx$
        avec $a_{n} = e^{2n}-1$, pour tout $n\in\N$.
        \begin{enumerate}[label=\textbf{\alph*)}, leftmargin=1.4em]
          \item Montrer que $U_{n} = 4n+2$. \hfill\textup{(0,75 pt)}
          \item Déterminer la nature de la suite $\left(U_{n}\right)$ en
                précisant sa raison. \hfill\textup{(0,5 pt)}
        \end{enumerate}
\end{enumerate}

% ---------------------------------------------------------------------------
\section*{L'état des sessions disponibles}

Contrairement à la série S, dont la banque d'annales s'étend sur plusieurs
sessions malgaches et sur de nombreux pays francophones, la série OSE dispose
d'un fonds propre beaucoup plus restreint. Quatre sessions malgaches ont été
repérées — 2021, 2022 et 2023 auprès de la médiathèque EDUCMAD, et 2026 dans
les archives locales de sujets officiels — mais seule la session 2026 a pu
être transcrite intégralement au moment de la rédaction de ce livre : c'est
elle, et elle seule, qui alimente les sections « annales » et « comme au
baccalauréat » des chapitres concernés.

\begin{center}
\setlength{\tabcolsep}{7pt}
\renewcommand{\arraystretch}{1.3}
\begin{tabular}{@{}L{4.3cm} L{8.4cm}@{}}
\toprule
\textsf{\footnotesize\bfseries EXERCICE DE LA SESSION 2026} &
\textsf{\footnotesize\bfseries UTILISÉ AU CHAPITRE} \\
\midrule
Exercice 1 — urne, variable aléatoire, loi binomiale &
  14, comme au baccalauréat \\
Exercice 2 — statistique à deux variables &
  15, seule annale sourcée \\
Exercice 3, parties I et II — intérêt simple et composé &
  annexe, volet « Intérêts et escompte », seule annale sourcée \\
Exercice 3, partie III — annuités constantes &
  annexe, volet « Annuités et emprunts », seule annale sourcée \\
Problème — étude de fonction, réciproque, suite intégrale &
  6, comme au baccalauréat \\
\bottomrule
\end{tabular}
\end{center}

\noindent
Les chapitres 1 à 5, 7 à 13, et les quatre chapitres de géométrie (9 à 12),
s'appuient sur un fonds d'annales générique déjà vérifié pour le livre de
série S — sujets de plusieurs pays francophones et sessions malgaches
d'autres séries —, ce contenu n'étant testé, à ce jour, dans aucun sujet OSE
identifié. Là où la banque malgache OSE s'est révélée trop mince pour fournir
un exercice « je cherche » ou « comme au baccalauréat » distinct de
l'annale — chapitre 15 et les deux volets de l'annexe financière —, ce
livre compose un exercice original,
au format de l'épreuve, plutôt que de dupliquer l'unique sujet réel
disponible. Chaque exercice de ce type porte la mention explicite
« exercice original ».

% ---------------------------------------------------------------------------
\section*{Conseils pour les quatre heures}

Le sujet n'impose aucun ordre. Commencez par l'exercice que vous savez
traiter : les points pris tôt libèrent l'esprit pour le reste.

Réservez au moins une heure et demie au problème d'analyse : il vaut neuf
points, près de la moitié de l'épreuve, et ses premières questions —
variations de la fonction auxiliaire, limites, tableau de variation — sont
les plus accessibles de tout le sujet. Ne les laissez pas de côté pour vous
lancer trop tôt dans la réciproque ou la suite intégrale, qui viennent en
fin de problème et rapportent moins de points pour un temps de rédaction
souvent plus long.

Dans l'exercice de mathématiques financières, identifiez d'abord ce qui est
donné — capital, taux, durée, valeur acquise — avant de choisir la formule :
la moitié des erreurs vient d'un intérêt simple confondu avec un intérêt
composé, ou d'une annuité confondue avec un intérêt.

Rédigez les justifications attendues, même brèves. « $f$ est continue et
strictement croissante sur $\intervalle{a}{b}$, et $k$ est compris entre
$f(a)$ et $f(b)$ » vaut le point ; le résultat seul ne le vaut pas.

Enfin, gardez dix minutes pour relire les signes et les unités — Ariary,
pourcentages, tonnes. C'est là, plus que dans les idées, que se jouent deux
ou trois points chaque année.
